\paragraph{What the diagnostics reveal.}
Our experiments suggest a practical gap between \emph{learning on a graph} and \emph{learning by enumerating candidates}.
Even when ground-truth trajectories are well aligned to the road graph, $K$-shortest candidate sets can miss many realized corridors, making candidate classification a weak foundation for modeling route distributions.

\paragraph{Why node-level AR is hard.}
Autoregressive generation over road-graph neighbors avoids geometric averaging, but it introduces a different brittleness: compounding error over hundreds (or thousands) of node transitions.
This is not merely an optimization issue---it is a horizon issue.
It motivates moving the corridor decision to a sparse representation (waypoints) so that the generative horizon matches the semantic decision horizon.

\paragraph{Waypoint AR + A* as a commitment--execution interface.}
Waypoint commitments are a natural interface between learning and planning.
The decision model can focus on selecting a small number of corridor-defining nodes conditioned on context, while the executor guarantees feasibility and connectivity on the graph.
Our current prototype highlights the remaining bottleneck: a coarse waypoint binning can produce unreachable waypoint pairs.
This points to reachability-aware decoding (e.g., distance-to-destination features, constrained sampling, or search-in-the-loop) as a necessary ingredient for a robust graph cascade.

\paragraph{Limitations and scope.}
We currently evaluate corridor match in graph space; generating realistic continuous trajectories within corridors is an important extension.
Map-matching and road-graph quality can also affect training labels and evaluation.
Finally, corridor-conditioned choice signal is data-hungry: cities with many repeated OD trips (e.g., Columbus) provide stronger evidence for context-conditioned diversity than smaller subsets (e.g., Detroit core).
